% ============================================================================
%  D/18-an-du-va-khung-tu-duy — file chương
%  Ngày tạo: 2026-09-07 | Sửa gần nhất: 2026-09-07 | Ôn gần nhất: chưa
%  Dependency: A/01, B/08, D/16. Mức độ: cốt lõi.
% ============================================================================
\documentclass[../../english.tex]{subfiles}
\begin{document}
\chapter{Ẩn dụ và khung tư duy (Metaphor and Framing)}
\ptrtarget{D/18}\ptrtarget{D/18-an-du-va-khung-tu-duy}
\dependency{\ptr{A/01} cho nền ``nghĩa nằm ở cách dùng''; \ptr{B/08} cho việc đọc ra thái độ; \ptr{D/16} cho màu sắc đánh giá của từ.}

\begin{tomtat}
Chương cuối của phần D nói về mức lớn nhất mà ngôn ngữ định hình hiểu biết: ẩn dụ. Ý trung tâm gây bất ngờ với nhiều người --- ẩn dụ không phải trang trí văn chương mà là cơ chế cơ bản để nghĩ về những thứ trừu tượng, và mọi ngành khoa học đều vận hành trên một hệ ẩn dụ mà người trong ngành thường không nhận ra. Nội dung gồm: ẩn dụ ý niệm và cách nhận ra chúng; ẩn dụ chết trong từ vựng học thuật; ẩn dụ nền trong các ngành kỹ thuật và cái chúng che khuất; định khung và cách nó hoạt động khi bạn đọc; và cách dùng ẩn dụ có trách nhiệm khi viết. Nếu chỉ nhớ một điều: mỗi ẩn dụ vừa cho bạn thấy một số thứ vừa \emph{che} một số thứ khác, và chỗ bị che là chỗ những câu hỏi chưa ai đặt ra đang nằm.
\end{tomtat}

\begin{nentang}
\kn{ẩn dụ ý niệm}{việc hiểu một miền trừu tượng qua một miền cụ thể: thời gian như tiền bạc, tranh luận như chiến tranh.}
\kn{miền nguồn}{miền cụ thể được mượn để nói: chiến tranh, dòng chảy, xây dựng.}
\kn{miền đích}{miền trừu tượng cần hiểu: tranh luận, thông tin, lý thuyết.}
\kn{ẩn dụ chết}{ẩn dụ đã dùng nhiều tới mức không ai còn thấy nó là ẩn dụ: \emph{grasp an idea}, \emph{a field of research}.}
\kn{định khung}{chọn cách mô tả khiến người nghe nghĩ theo một hướng nhất định.}
\kn{che khuất}{mặt của hiện tượng mà một ẩn dụ làm người ta không nhìn tới.}
\end{nentang}

\begin{khoigoi}
\h{Bạn có dùng ẩn dụ khi viết bài kỹ thuật không? Nếu nghĩ là không, hãy thử giải thích cụm \emph{a strong argument}.}
\h{Ẩn dụ giúp hiểu --- nhưng nó có mất mát gì không, và mất ở đâu?}
\h{Hai bản tin cùng số liệu, một bản nói ``thất nghiệp tăng lên 6\%'', một bản nói ``95\% người dân vẫn có việc''. Cơ chế nào đang hoạt động?}
\h{Ngành của bạn dựa trên ẩn dụ nào? Làm sao tìm ra nó?}
\h{Khi viết, dùng ẩn dụ thế nào là có trách nhiệm?}
\end{khoigoi}

Hãy bắt đầu bằng câu hỏi mở đầu thứ nhất. Nhiều người viết kỹ thuật tin rằng bài của họ không có ẩn dụ nào --- ẩn dụ là chuyện của thơ. Nhưng hãy nhìn vào các cụm bạn dùng hằng ngày.

Bạn \emph{build} một mô hình. Bạn \emph{grasp} một khái niệm. Một lý thuyết có \emph{foundations} và có thể \emph{collapse}. Một lập luận \emph{strong} hoặc \emph{weak}, nó có thể bị \emph{attacked} và cần được \emph{defended}. Thông tin \emph{flows}. Một ngành là một \emph{field} mà bạn \emph{explore}.

Mỗi cụm trong số đó mượn một miền cụ thể để nói về một thứ trừu tượng. Lý thuyết không thật sự có móng; lập luận không thật sự bị tấn công. Đây là ẩn dụ --- chỉ là loại đã dùng nhiều tới mức bạn không còn thấy nó nữa.

\section{Ẩn dụ không phải trang trí}

Quan sát nền của Lakoff và Johnson, trình bày năm 1980\cite{lakoff1980}, là: các ẩn dụ này không rời rạc mà tụ thành \emph{hệ}. Nếu bạn thu thập mọi cách người ta nói về tranh luận, bạn sẽ thấy chúng đến từ một miền nguồn duy nhất.

\begin{center}\footnotesize
\rowcolors{2}{white}{tblalt}
\begin{longtable}{@{}L{4.4cm}L{4.6cm}L{5.4cm}@{}}
\toprule
\textbf{Ẩn dụ nền} & \textbf{Cách nói kéo theo} & \textbf{Cái nó che khuất}\\
\midrule
\endfirsthead
\toprule
\textbf{Ẩn dụ nền} & \textbf{Cách nói kéo theo} & \textbf{Cái nó che khuất (tiếp)}\\
\midrule
\endhead
Tranh luận là chiến tranh & \emph{attack a claim}, \emph{defend a position}, \emph{win an argument} & Tranh luận cũng có thể là hợp tác tìm ra sự thật\\
Lý thuyết là toà nhà & \emph{foundations}, \emph{framework}, \emph{collapse}, \emph{support} & Lý thuyết có thể tiến hoá, không chỉ đứng hay đổ\\
Thời gian là tiền & \emph{spend time}, \emph{waste time}, \emph{save time} & Thời gian không tích luỹ được như tiền\\
Ý tưởng là vật thể & \emph{grasp}, \emph{hold}, \emph{transfer}, \emph{get across} & Người nhận \emph{xây lại} ý chứ không nhận nguyên khối\\
Thông tin là chất lỏng & \emph{flow}, \emph{leak}, \emph{stream}, \emph{pour in} & Thông tin được diễn giải, không chảy thụ động\\
Hiểu là nhìn & \emph{I see}, \emph{clear}, \emph{illuminate}, \emph{obscure} & Hiểu cũng có thể là làm được, không chỉ thấy được\\
\bottomrule
\end{longtable}
\end{center}

Cột thứ ba là phần quan trọng nhất của chương này và là câu trả lời cho câu hỏi mở đầu thứ hai. Mỗi ẩn dụ làm nổi bật một số mặt của hiện tượng và \emph{che khuất} các mặt khác. Nếu bạn chỉ có ẩn dụ chiến tranh cho tranh luận, thì việc thay đổi ý kiến sau khi nghe người khác sẽ được cảm nhận là \emph{thua} --- một cách hiểu tai hại cho khoa học, và nó đến từ ngôn ngữ chứ không từ logic.

\begin{figure}[H]\centering
\begin{tikzpicture}[
  src/.style={draw=steel,fill=bluebg,rounded corners=3pt,inner sep=6pt,align=center,font=\small,text width=3.0cm},
  tgt/.style={draw=greendark,fill=greenbg,rounded corners=3pt,inner sep=6pt,align=center,font=\small,text width=3.0cm},
  hid/.style={draw=reddark,dashed,fill=redbg,rounded corners=3pt,inner sep=5pt,align=center,font=\scriptsize,text width=3.2cm},
  ar/.style={-{Latex[length=2.4mm]},draw=accent,thick},
  lb/.style={font=\scriptsize\itshape,color=tiercol,align=center}]
\node[src] (s) {\textbf{Miền nguồn}\\ Chiến tranh\\ {\scriptsize tấn công, phòng thủ, thắng thua}};
\node[tgt,right=3.0cm of s] (t) {\textbf{Miền đích}\\ Tranh luận};
\draw[ar] (s) -- node[lb,above=1pt,text width=2.6cm] {ánh xạ: cấu trúc của nguồn áp lên đích} (t);
\node[hid,below=0.9cm of t] (h) {\textbf{Bị che:} tranh luận như hợp tác; đổi ý là tiến bộ, không phải thua};
\draw[-{Latex[length=2.0mm]},draw=reddark,dashed,thick] (t) -- (h);
\end{tikzpicture}
\caption{Cấu trúc của một ẩn dụ ý niệm. Ánh xạ mang cấu trúc của miền nguồn sang miền đích --- và cái không có trong miền nguồn thì trở nên khó nghĩ tới trong miền đích.}
\label{fig:an-du-anh-xa}
\end{figure}

\section{Ẩn dụ nền trong các ngành kỹ thuật}

Câu hỏi mở đầu thứ tư hỏi cách tìm ẩn dụ của ngành bạn. Có một cách rất cụ thể: \emph{tìm các động từ}. Danh từ trong một ngành thường là thuật ngữ trung tính, nhưng động từ đi với chúng lại lộ ra miền nguồn.

Vài ví dụ từ các ngành kỹ thuật, và cái chúng che.

\emph{Bộ nhớ máy tính là một kho chứa.} Ta \emph{store}, \emph{retrieve}, \emph{fetch}, và nói về \emph{locations}. Ẩn dụ này rất tốt cho bộ nhớ máy tính. Nhưng khi cùng ẩn dụ đó được dùng cho \emph{trí nhớ con người}, nó che mất một điều mà nghiên cứu tâm lý đã chỉ ra rõ\cite{roediger2006}: nhớ lại là một quá trình \emph{tái dựng}, không phải lấy ra một bản lưu.

\emph{Mạng thần kinh nhân tạo là một bộ não.} Ẩn dụ này giúp trực giác ban đầu, nhưng nó che mất khoảng cách rất lớn giữa hai thứ và làm nảy sinh những suy diễn không có cơ sở về việc mô hình ``hiểu'' hay ``suy nghĩ''.

\emph{Mã nguồn là một toà nhà.} Ta nói \emph{build}, \emph{architecture}, \emph{foundation}, \emph{technical debt}. Ẩn dụ này che mất một điểm quan trọng: phần mềm được sửa liên tục, còn toà nhà thì không --- và nhiều cách quản lý dự án hỏng vì áp thẳng trực giác xây dựng lên phần mềm.

\emph{Ước lượng là theo dõi một mục tiêu.} Trong các hệ định vị và cảm biến, ta \emph{track}, \emph{lose}, \emph{recover}, \emph{drift}. Ẩn dụ săn đuổi này rất tự nhiên và cũng hữu ích, nhưng nó có thể che mờ việc đối tượng thật sự được ước lượng là một \emph{phân bố xác suất} chứ không phải một điểm đang chạy.

Điểm cần rút: không ẩn dụ nào ``sai''. Chúng là công cụ tư duy, và mỗi công cụ có vùng hoạt động tốt. Cái nguy hiểm là quên rằng mình đang dùng một ẩn dụ, và suy ra các kết luận từ miền nguồn mà miền đích không chống đỡ.

\begin{cannho}
Một kỹ thuật nghiên cứu dùng được ngay: khi bí ở một vấn đề, hãy hỏi ``mình đang dùng ẩn dụ nào, và nếu đổi sang một ẩn dụ khác thì câu hỏi sẽ khác thế nào?''. Đổi từ ``theo dõi một điểm'' sang ``duy trì một niềm tin'' đặt ra những câu hỏi rất khác. Đây là một trong số ít cách có hệ thống để tìm những câu hỏi mà chưa ai trong ngành đặt ra --- vì chúng nằm đúng ở chỗ ẩn dụ hiện hành che khuất.
\end{cannho}

\section{Định khung}

Câu hỏi mở đầu thứ ba mô tả hiện tượng định khung. Hai bản tin cùng số liệu, nhưng một bản đặt trọng tâm ở phần mất và một bản đặt ở phần còn --- và người đọc phản ứng khác nhau, dù thông tin y hệt.

Cơ chế: một cách mô tả kích hoạt một khung, và khung mang theo cả một bộ liên tưởng, giá trị, và câu hỏi mặc định. Khi bạn nói \emph{tax relief}, bạn đã ngầm đặt thuế vào vai trò một gánh nặng cần được giảm nhẹ --- và người phản đối phải chống lại cả cái khung đó chứ không chỉ chống lại một chính sách.

Trong văn khoa học, định khung xuất hiện ở những chỗ tinh vi hơn nhưng cùng cơ chế.

\begin{tuvung}[Cùng sự việc, khác khung]
\tv{fails in 8\% of cases}{Hỏng trong 8\% trường hợp}{Khung thất bại; chú ý dồn vào phần hỏng}{Đúng nhưng dẫn người đọc tới câu hỏi ``vì sao hỏng''}
\tv{succeeds in 92\% of cases}{Thành công 92\%}{Khung thành công; cùng số liệu}{Cùng dữ liệu, phản ứng khác}
\tv{a limitation of the method}{Một hạn chế}{Khung khiếm khuyết}{Đặt trách nhiệm ở phương pháp}
\tv{a condition for validity}{Một điều kiện áp dụng}{Khung phạm vi}{Trung tính hơn, và thường chính xác hơn (\ptr{C/14})}
\tv{the algorithm decides}{Thuật toán quyết định}{Khung tác nhân có ý chí}{Che mất việc con người đã chọn tiêu chí}
\tv{the algorithm computes a score}{Thuật toán tính điểm}{Khung công cụ}{Giữ trách nhiệm ở người thiết kế}
\end{tuvung}

Hai dòng cuối đáng chú ý vì chúng có hệ quả ngoài ngôn ngữ. Viết \emph{the algorithm decides who gets a loan} khác hẳn \emph{the bank uses an algorithm to decide who gets a loan}. Câu đầu làm người đọc quên rằng có người đã chọn dữ liệu, chọn tiêu chí, và chịu trách nhiệm. Đây là một chỗ mà lựa chọn ngôn ngữ có hệ quả đạo đức thật, và nó nối với \ptr{B/08}: khi đọc, hãy để ý ai bị làm cho biến mất khỏi câu.

\section{Ẩn dụ chết và thuật ngữ}

Có một hiện tượng đáng chú ý ở ranh giới giữa chương này và \ptr{B/10}: rất nhiều thuật ngữ kỹ thuật \emph{là} ẩn dụ đã chết.

\emph{Kernel} là hạt nhân. \emph{Tree}, \emph{root}, \emph{leaf}, \emph{branch}. \emph{Bug}. \emph{Noise}. \emph{Bottleneck}. \emph{Pipeline}. \emph{Memory}. \emph{Virus}. \emph{Cloud}. Mỗi từ này từng là một so sánh sống động, và giờ chỉ là tên.

Điều này có hai hệ quả thực hành. Thứ nhất, khi học một thuật ngữ mới, hỏi ``đây là ẩn dụ từ đâu?'' thường giúp nhớ và hiểu nhanh hơn, vì nó nối khái niệm mới vào một thứ bạn đã biết --- một dạng của \ptr{A/04}.

Thứ hai, và quan trọng hơn: ẩn dụ chết vẫn tiếp tục \emph{định hướng suy nghĩ} dù không ai còn nhận ra. Gọi một lỗi phần mềm là \emph{bug} ngầm gợi ý nó là một thứ nhỏ, ngẫu nhiên, xâm nhập từ bên ngoài --- trong khi phần lớn lỗi là hệ quả trực tiếp của một quyết định thiết kế. Cách gọi ảnh hưởng tới cách người ta phản ứng.

\section{Dùng ẩn dụ có trách nhiệm khi viết}

Câu hỏi mở đầu thứ năm. Ẩn dụ là công cụ giải thích mạnh nhất bạn có, nhất là khi viết cho người ngoài ngành. Bốn nguyên tắc để dùng nó mà không gây hiểu sai.

\emph{Một --- nói rõ giới hạn ngay sau khi đưa ra.} \emph{The map behaves like a rubber sheet, in the sense that\ldots} Cụm \emph{in the sense that} làm rất nhiều việc: nó cho biết ánh xạ đi tới đâu.

\emph{Hai --- kiểm xem người đọc có thể suy diễn sai gì.} Với mỗi ẩn dụ, hỏi: từ miền nguồn, người đọc có thể rút ra kết luận nào mà miền đích không chống đỡ? Nếu có kết luận nguy hiểm, hãy chặn trước.

\emph{Ba --- tránh ẩn dụ tác nhân cho hệ thống tự động.} Như mục 3 đã chỉ ra, \emph{the model thinks}, \emph{the algorithm decides} có hệ quả ngoài ngôn ngữ. Dùng động từ mô tả cái hệ thống thật sự làm.

\emph{Bốn --- đừng trộn ẩn dụ.} \emph{We build a framework that flows through the pipeline} trộn ba miền nguồn và làm người đọc mất chỗ bám. Một ẩn dụ, giữ nhất quán trong cả đoạn.

\section{Gốc rễ: từ phép tu từ tới cơ chế tư duy}

Trong truyền thống tu từ học phương Tây, kéo dài từ Aristotle tới tận thế kỷ 20, ẩn dụ được xếp vào loại \emph{trang trí}: một cách nói bóng bẩy thay cho cách nói thẳng, thuộc về văn chương và diễn thuyết chứ không thuộc về tư duy nghiêm túc. Theo cách nhìn đó, một văn bản khoa học lý tưởng sẽ không có ẩn dụ nào.

Công trình của Lakoff và Johnson năm 1980\cite{lakoff1980} đảo ngược cách nhìn này, và lập luận của họ dựa trên bằng chứng ngôn ngữ chứ không trên tư biện: họ chỉ ra rằng các cách nói ẩn dụ trong đời thường không rời rạc mà tụ thành hệ có cấu trúc. Nếu ẩn dụ chỉ là trang trí, không có lý do gì hàng chục cách nói khác nhau về tranh luận lại cùng đến từ miền chiến tranh. Sự có hệ thống đó gợi ý rằng ẩn dụ nằm ở tầng \emph{ý niệm} chứ không ở tầng diễn đạt: ta không chỉ nói về tranh luận bằng từ ngữ chiến tranh, ta \emph{nghĩ} về nó theo cách đó.

Ý này gặp một dòng nghiên cứu song song về \emph{định khung} trong tâm lý học ra quyết định. Các nghiên cứu ở đó cho thấy cùng một lựa chọn được mô tả theo hướng được và theo hướng mất dẫn tới quyết định khác nhau ở cùng những người --- một kết quả khó dung hoà với giả định rằng cách diễn đạt chỉ là lớp vỏ.

Trong lịch sử khoa học, vai trò của ẩn dụ còn rõ hơn nữa: nhiều bước tiến lớn gắn với một ẩn dụ mới chứ không chỉ với dữ liệu mới. Nghĩ về ánh sáng như sóng, về nhiễm bệnh như một cuộc xâm nhập của sinh vật nhỏ, về di truyền như một mã --- mỗi ẩn dụ mở ra một tập câu hỏi mà ẩn dụ trước không cho phép đặt. Điều này giải thích vì sao kỹ thuật ở mục 2 --- đổi ẩn dụ để tìm câu hỏi mới --- không phải một trò chơi ngôn ngữ mà là một phương pháp có tiền lệ.

Cuối cùng, một cảnh báo mà chính dòng nghiên cứu này đưa ra: vì ẩn dụ định hình cái ta nghĩ tới được, việc một cộng đồng chỉ có \emph{một} ẩn dụ cho đối tượng của mình là một rủi ro nhận thức. Đó là lý do cột ``cái nó che khuất'' ở mục 1 đáng để bạn tự điền cho ngành của mình.

\begin{gocre}
\gr{Aristotle tới thế kỷ 20}{Ẩn dụ là gì?}{Xếp vào phép trang trí; văn bản khoa học lý tưởng thì không có ẩn dụ.}
\gr{1980 --- Lakoff \& Johnson\cite{lakoff1980}}{Vì sao các cách nói ẩn dụ lại có hệ thống?}{Ẩn dụ nằm ở tầng ý niệm, không ở tầng diễn đạt --- ta nghĩ bằng ẩn dụ.}
\gr{Tâm lý học ra quyết định}{Cách mô tả có ảnh hưởng tới lựa chọn không?}{Định khung được--mất đổi quyết định ở cùng những người \(\Rightarrow\) diễn đạt không phải lớp vỏ.}
\gr{Lịch sử khoa học}{Bước tiến đến từ đâu?}{Nhiều bước tiến gắn với một ẩn dụ mới (sóng, mầm bệnh, mã di truyền) chứ không chỉ dữ liệu mới.}
\gr{Ẩn dụ chết thành thuật ngữ}{Thuật ngữ kỹ thuật từ đâu ra?}{\emph{kernel}, \emph{tree}, \emph{bug}, \emph{noise} --- ẩn dụ nguội đi thành tên gọi, nhưng vẫn định hướng suy nghĩ.}
\gr{Cảnh báo về đơn ẩn dụ}{Một cộng đồng chỉ có một ẩn dụ thì sao?}{Rủi ro nhận thức: các câu hỏi bị che khuất không ai đặt ra.}
\end{gocre}

\section{Liên hệ ngang}

\begin{lienhe}
\lh{Lịch sử khoa học}{Đổi mô hình nền}{Nhiều cuộc đổi mô hình đi kèm đổi ẩn dụ; kỹ thuật ở mục 2 là phiên bản cá nhân của cùng cơ chế.}
\lh{Tâm lý học}{Trí nhớ là tái dựng, không phải lấy ra}{Ví dụ rõ nhất của việc ẩn dụ kho chứa che mất bản chất hiện tượng.}
\lh{Chính trị và truyền thông}{Chọn từ để định khung}{Cùng cơ chế mục 3 ở quy mô lớn; đọc ra nó là một phần của \ptr{B/08}.}
\lh{Đạo đức công nghệ}{Ai chịu trách nhiệm khi hệ thống tự động sai}{Ẩn dụ tác nhân (\emph{the algorithm decides}) làm biến mất người ra quyết định --- hệ quả pháp lý thật.}
\lh{Toán học}{Ký hiệu tốt gợi ý hướng đi}{Ký hiệu cũng là một dạng ẩn dụ đóng gói; ký hiệu tốt làm một số phép biến đổi trở nên hiển nhiên.}
\lh{Dạy học}{Giải thích bằng so sánh}{Ẩn dụ là công cụ dạy mạnh nhất và cũng nguy hiểm nhất --- người học giữ lại cả phần suy diễn sai nếu bạn không chặn trước (mục 5).}
\end{lienhe}

\begin{traloi}
\tl{Có, rất nhiều. \emph{A strong argument} mượn miền sức mạnh vật lý; lập luận không thật sự có sức. Tương tự: bạn \emph{build} một mô hình, \emph{grasp} một khái niệm, một lý thuyết có \emph{foundations} và có thể \emph{collapse}, thông tin \emph{flows}, một ngành là một \emph{field} bạn \emph{explore}. Đây là ẩn dụ đã dùng nhiều tới mức không ai còn thấy --- ẩn dụ chết --- nhưng chúng vẫn định hướng suy nghĩ.}
\tl{Có, và mất ở chỗ \emph{che khuất}. Mỗi ẩn dụ làm nổi bật một số mặt và làm các mặt khác trở nên khó nghĩ tới. Nếu tranh luận chỉ là chiến tranh thì đổi ý sau khi nghe người khác được cảm nhận là thua --- một cách hiểu tai hại cho khoa học, đến từ ngôn ngữ chứ không từ logic. Nếu bộ nhớ chỉ là kho chứa thì việc nhớ lại là tái dựng bị che mất. Chỗ bị che là chỗ những câu hỏi chưa ai đặt ra đang nằm.}
\tl{Định khung: cùng dữ liệu nhưng cách mô tả kích hoạt một khung mang theo cả bộ liên tưởng, giá trị và câu hỏi mặc định. Bản nói ``thất nghiệp tăng lên 6\%'' dẫn chú ý tới phần mất; bản nói ``95\% vẫn có việc'' dẫn tới phần còn. Trong văn khoa học cùng cơ chế xuất hiện tinh vi hơn: \emph{a limitation of the method} (khung khiếm khuyết) so với \emph{a condition for validity} (khung phạm vi) --- cách sau thường chính xác hơn.}
\tl{Cách tìm cụ thể: nhìn vào các \emph{động từ}, vì danh từ trong ngành thường là thuật ngữ trung tính còn động từ lộ ra miền nguồn. Ví dụ: \emph{store}, \emph{retrieve}, \emph{fetch} cho thấy ẩn dụ kho chứa; \emph{build}, \emph{architecture}, \emph{foundation} cho thấy ẩn dụ toà nhà; \emph{track}, \emph{lose}, \emph{recover}, \emph{drift} cho thấy ẩn dụ săn đuổi. Rồi tự điền cột ``cái nó che khuất'' cho ngành mình --- đó là chỗ tìm được câu hỏi mới.}
\tl{Bốn nguyên tắc. Nói rõ giới hạn ngay sau khi đưa ra ẩn dụ --- cụm \emph{in the sense that} rất hiệu quả. Kiểm xem người đọc có thể rút ra kết luận sai nào từ miền nguồn, và chặn trước. Tránh ẩn dụ tác nhân cho hệ thống tự động (\emph{the algorithm decides}) vì nó làm biến mất người chịu trách nhiệm. Và đừng trộn ẩn dụ --- \emph{build a framework that flows through the pipeline} trộn ba miền nguồn và làm người đọc mất chỗ bám.}
\end{traloi}

\begin{dongian}
Nhiều người nghĩ ẩn dụ là chuyện của thơ, còn viết kỹ thuật thì không có. Thử kiểm lại xem.

Bạn \emph{xây} một mô hình. Bạn \emph{nắm} được một khái niệm. Một lý thuyết có \emph{nền móng} và có thể \emph{sụp đổ}. Một lập luận thì \emph{mạnh} hoặc \emph{yếu}, và người ta \emph{tấn công} hoặc \emph{bảo vệ} nó. Thông tin thì \emph{chảy}.

Không cái nào trong số đó là nghĩa đen. Lý thuyết không có móng thật; lập luận không bị tấn công thật. Đó đều là ẩn dụ --- chỉ là loại quen tới mức trong suốt.

Điều đáng nói là chúng không vô hại. Mỗi ẩn dụ cho bạn thấy vài thứ và \emph{che} vài thứ. Nếu tranh luận trong đầu bạn là một cuộc chiến, thì đổi ý sau khi nghe người khác nói sẽ có cảm giác như thua --- trong khi với khoa học, đó chính là tiến bộ. Cảm giác đó đến từ ngôn ngữ, không từ thực tế.

Chuyện này rất thực trong ngành kỹ thuật. Ta nói bộ nhớ máy tính là cái kho: \emph{lưu vào}, \emph{lấy ra}. Ẩn dụ tốt cho máy. Nhưng khi mượn nó cho trí nhớ con người, nó che mất một điều đã được chứng minh: nhớ lại là \emph{dựng lại}, không phải lấy ra một bản đã lưu.

Từ đây có một mẹo dùng được cho công việc nghiên cứu: khi bí, hãy hỏi ``mình đang dùng hình ảnh nào để nghĩ về cái này, và nếu đổi sang hình ảnh khác thì câu hỏi sẽ khác thế nào?''. Chuyển từ ``đuổi theo một điểm'' sang ``duy trì một niềm tin'' làm nảy ra những câu hỏi hoàn toàn khác. Chỗ ẩn dụ cũ che khuất chính là chỗ chưa ai tìm.

Và một lưu ý khi viết: cẩn thận với những câu như \emph{thuật toán quyết định ai được vay}. Nghe thì gọn, nhưng nó làm biến mất những người đã chọn dữ liệu và chọn tiêu chí. Viết \emph{ngân hàng dùng một thuật toán để quyết định} dài hơn vài chữ nhưng nói đúng ai chịu trách nhiệm.
\motcau{Mỗi ẩn dụ cho thấy một phần và giấu đi phần còn lại.}
\chuadongian{Biết một ẩn dụ đang che mất cái gì đòi hiểu biết về chính hiện tượng đó --- nên phần này không tách khỏi chuyên môn được.}
\end{dongian}

\begin{onbai}
\ar[3]{Ẩn dụ không phải trang trí}{Nó là cơ chế cơ bản để nghĩ về những thứ trừu tượng.}
\ar[3]{Cấu trúc ẩn dụ ý niệm}{Miền nguồn cụ thể ánh xạ cấu trúc sang miền đích trừu tượng.}
\ar[3]{Cái quan trọng nhất}{Mỗi ẩn dụ vừa làm nổi bật vừa \emph{che khuất}; chỗ bị che là chỗ có câu hỏi chưa ai đặt.}
\ar[3]{Cách tìm ẩn dụ của ngành}{Nhìn vào các \emph{động từ}, không nhìn danh từ.}
\ar[2]{Ví dụ ẩn dụ nền}{Tranh luận là chiến tranh; lý thuyết là toà nhà; ý tưởng là vật thể; hiểu là nhìn.}
\ar[2]{Ẩn dụ kho chứa}{Tốt cho bộ nhớ máy; che mất việc nhớ lại ở người là tái dựng.}
\ar[2]{Định khung}{Cùng dữ liệu, cách mô tả kích hoạt khung khác nhau và đổi phản ứng của người đọc.}
\ar[2]{Khung khiếm khuyết và khung phạm vi}{\emph{a limitation} so với \emph{a condition for validity} --- cách sau thường chính xác hơn.}
\ar[2]{Ẩn dụ tác nhân}{\emph{The algorithm decides} làm biến mất người chịu trách nhiệm; có hệ quả đạo đức thật.}
\ar[2]{Ẩn dụ chết thành thuật ngữ}{\emph{kernel}, \emph{tree}, \emph{bug}, \emph{noise}, \emph{pipeline} --- vẫn định hướng suy nghĩ dù không ai nhận ra.}
\ar[2]{Bốn nguyên tắc dùng ẩn dụ}{Nói rõ giới hạn; chặn suy diễn sai; tránh ẩn dụ tác nhân; đừng trộn ẩn dụ.}
\ar[1]{Kỹ thuật tìm câu hỏi mới}{Đổi ẩn dụ và xem câu hỏi thay đổi thế nào --- có tiền lệ trong lịch sử khoa học.}
\end{onbai}

\begin{hoidap}
\qa{Trong bài báo tôi có nên dùng ẩn dụ không?}{Có, nhưng có chỗ nên và chỗ không. Nên trong phần mở bài và phần thảo luận, nơi bạn phải làm người đọc hình dung được ý tưởng; và đặc biệt nên khi giải thích cho người ngoài chuyên môn hẹp. Không nên trong phần phương pháp và kết quả, nơi cần mô tả chính xác cái bạn làm và đo. Và bất kể ở đâu, hãy áp nguyên tắc một ở mục 5: nói rõ ánh xạ đi tới đâu.}
\qa{Ẩn dụ có làm bài nghe kém nghiêm túc không?}{Không, nếu bạn dùng nó để giải thích chứ không để thay thế lập luận. Điều thực sự làm bài kém nghiêm túc là ẩn dụ \emph{không được giới hạn} --- câu kiểu ``mô hình học giống như một đứa trẻ'' để trần, không nói giống ở điểm nào. Ẩn dụ có kèm \emph{in the sense that} thì ngược lại: nó cho thấy bạn nắm rõ đối tượng đủ để biết so sánh này đúng tới đâu.}
\qa{Làm sao tôi biết ẩn dụ trong ngành mình đang che mất gì?}{Đây là câu hỏi khó nhất của chương, và không có cách máy móc. Hai gợi ý làm được. Thứ nhất, tìm những hiện tượng mà ngành bạn liên tục gọi là ``bất thường'' hoặc ``ngoại lệ'' --- chúng thường nằm đúng ở chỗ ẩn dụ nền không mô tả được. Thứ hai, xem một ngành lân cận mô tả cùng hiện tượng bằng ẩn dụ nào, rồi hỏi ẩn dụ của họ nhìn ra cái gì mà ẩn dụ của bạn không.}
\qa{Đọc ra định khung có làm tôi thành đa nghi quá không?}{Rủi ro có thật, và cách phòng là phân biệt hai câu hỏi. Câu hỏi thứ nhất: cách mô tả này có khung không? Câu trả lời gần như luôn là có --- mọi cách mô tả đều có khung, kể cả của bạn, vì không tồn tại cách mô tả trung tính tuyệt đối. Câu hỏi thứ hai, mới là câu đáng hỏi: khung này có làm người đọc rút ra kết luận mà dữ liệu không chống đỡ không? Chỉ khi trả lời có thì mới đáng nêu ra.}
\qa{Ẩn dụ có khác nhau giữa các ngôn ngữ không?}{Có và không, và cả hai đều thú vị. Một số ẩn dụ rất phổ quát --- nhiều thì gắn với lên trên, quan trọng gắn với lớn, thời gian gắn với chuyển động --- có lẽ vì chúng bắt nguồn từ trải nghiệm thân thể chung. Nhưng nhiều ẩn dụ khác thì đặc thù văn hoá, và đây là nguồn hiểu sai khi dịch. Khi bạn dịch một ẩn dụ từ tiếng Việt sang tiếng Anh, hãy áp phép thử ba câu ở \ptr{D/16} --- nếu người ta không nói thế, hãy diễn đạt thẳng thay vì giữ hình ảnh.}
\qa{Có ẩn dụ nào trong ngôn ngữ về việc học không, và nó có che gì không?}{Có, và nó đáng để bạn nhìn kỹ vì nó ảnh hưởng tới chính cách bạn học. Ẩn dụ phổ biến nhất là học như \emph{thu nhận}: bạn \emph{lấy} kiến thức, \emph{nạp} thông tin, \emph{tích luỹ} vốn từ. Nó che mất một mặt quan trọng: học cũng là \emph{tham gia} vào một cộng đồng thực hành --- bạn học tiếng Anh khoa học không phải bằng cách chứa thêm từ mà bằng cách làm những việc mà cộng đồng đó làm. Phần \ptr{E/19} xây trên cách nhìn thứ hai\cite{nation2013}.}
\qa{Ẩn dụ chết có nên hồi sinh không khi viết?}{Đôi khi rất hiệu quả, và đây là một kỹ thuật đáng biết. Khi bạn muốn người đọc suy nghĩ lại về một khái niệm quen, hãy kéo ẩn dụ chết của nó sống lại: nhắc rằng \emph{bug} vốn có nghĩa là con bọ, rồi hỏi liệu ẩn dụ đó có đúng cho phần lớn lỗi phần mềm không. Kỹ thuật này mạnh nên dùng tiết chế --- một lần trong một bài là đủ, và chỉ khi cái bạn muốn nói thật sự nằm ở chỗ ẩn dụ che khuất.}
\end{hoidap}

\begin{baitap}
\bt[1]{Tìm 10 ẩn dụ chết trong một bài báo ngành bạn}{Bài báo ngành}{Bắt đầu từ các động từ, không từ danh từ.}
\bt[1]{Với 5 thuật ngữ trong ngành, truy xem nó là ẩn dụ từ miền nào}{Bảng thuật ngữ của bạn}{Ghi vào bảng ở \ptr{B/10} --- giúp nhớ và hiểu.}
\bt[2]{Điền cột ``cái nó che khuất'' cho 3 ẩn dụ nền của ngành bạn}{Tài liệu ngành}{Đây là bài khó nhất và cũng đáng giá nhất của chương.}
\bt[2]{Viết lại 5 câu có ẩn dụ tác nhân thành câu nêu rõ người chịu trách nhiệm}{Bài viết của bạn}{\emph{The algorithm decides} \(\rightarrow\) \emph{We use an algorithm to \ldots}}
\bt[2]{Tìm 3 cặp câu cùng dữ liệu khác khung trong tài liệu hoặc tin tức}{Tự chọn}{Ghi rõ khung nào dẫn tới câu hỏi mặc định nào.}
\bt[3]{Chọn một vấn đề bạn đang bí và thử ba ẩn dụ khác nhau cho nó}{Công việc của bạn}{Ghi lại câu hỏi mới nảy ra từ mỗi ẩn dụ.}
\bt[3]{Viết một đoạn giải thích một khái niệm ngành bạn cho người ngoài, dùng một ẩn dụ có giới hạn rõ}{Bài viết của bạn}{Bắt buộc có cụm \emph{in the sense that} hoặc tương đương.}
\end{baitap}

\section*{Kết chương và đọc tiếp}

Phần D khép lại ở đây. Ba chương đi qua ba mức mà nghĩa được quyết định ngoài định nghĩa từ điển: sắc thái và văn vực ở mức từ, thói quen kết hợp ở mức tổ hợp, và ẩn dụ ở mức hệ khái niệm. Cùng nhau chúng trả lời mục tiêu bạn đặt ra ban đầu --- hiểu sâu về ngôn từ --- theo nghĩa cụ thể nhất: biết cái gì đang quyết định nghĩa mà không nằm trong định nghĩa.

Bốn phần đầu đã cho bạn hiểu biết. Nhưng hiểu biết về ngôn ngữ không tự chuyển thành khả năng dùng, và khoảng cách giữa hai thứ đó lớn hơn ở ngôn ngữ so với hầu hết các lĩnh vực khác --- đó là lý do có rất nhiều người biết nhiều về tiếng Anh mà dùng không trôi. Phần E xử lý đúng khoảng cách đó: nó nói về việc luyện hằng ngày, về công cụ, và về một lộ trình cụ thể để biến những gì bạn vừa đọc thành thói quen. Bắt đầu từ \ptr{E/19}.
\ifSubfilesClassLoaded{\printbibliography[title={Nguồn}]}{}
\end{document}
